\NativeSlot{45}{90}{%
\ExamGroup{三、计算。{\normalfont\kaishu\fontsize{10.8bp}{14bp}\selectfont （共24分）}}

\noindent \textbf{1.}\quad 列竖式计算。（8分）\par\smallskip
\noindent%
% p015_s01
\makebox[.249\linewidth][l]{\(258\times13=\hspace{1.8em}\)}%
% p015_s02
\makebox[.249\linewidth][l]{\(754\div26=\hspace{1.8em}\)}%
% p015_s03
\makebox[.249\linewidth][l]{\(308\times25=\hspace{1.8em}\)}%
% p015_s04
\makebox[.249\linewidth][l]{\(552\div69=\hspace{1.8em}\)}\par
}
\NativeSlot{197}{64}{%
\noindent \textbf{2.}\quad 计算下列各题。（16分）\par\smallskip
\par\noindent%
\ExamTile{.485}{%
% p015_s05

\noindent \(4000-15\times(24+195)\)\par

\ExamWork{3}
}\hfill%
\ExamTile{.485}{%
% p015_s06

\noindent \(180\div[(84-12)\div4]\)\par

\ExamWork{3}
}\par
}
\NativeSlot{324}{53}{%
\par\noindent%
\ExamTile{.485}{%
% p015_s07

\noindent \(34\times[160-(60-15)]\)\par

\ExamWork{3}
}\hfill%
\ExamTile{.485}{%
% p015_s08

\noindent \([288\div(295-289)]\times17\)\par

\ExamWork{3}
}\par
}
\NativeSlot{423}{100}{%
\ExamGroup{四、画一画。{\normalfont\kaishu\fontsize{10.8bp}{14bp}\selectfont （共5分）}}

% p015_s09

\ExamQuestion{画一画}{5}%
过点\(A\)作已知直线的垂线段，在该垂线段上取点\(B\)，使\(AB=3\)厘米，再过点\(B\)画出已知直线的平行线。\par
\BookFigure{figures/p015_s09-01.tex}
}
\NativeSlot{599}{117}{%
\ExamGroup{五、解决问题。{\normalfont\kaishu\fontsize{10.8bp}{14bp}\selectfont （共30分）}}

% p015_s10

\ExamQuestion{1.}{5}%
小鸭很擅长游泳，在陆地上走路却很费劲。它打算从自己家到小鸡家玩耍（如图）。想一想，怎样走才能最省力气？请画出路线图。\par
\BookFigure{figures/p015_s10-01.tex}
}
